\chapter{Base change and permanent models}\label{ch:base-change-and-permanent-models}
\chaptermark{Base change and permanent models}

This chapter studies the behaviour of models and their component valuations under finite extensions of the base field. Its main themes are normalized base change, reduced special fibres and permanent models, the elimination of ramification, and illustrative tame and wild examples.

Throughout this chapter, $X$ is a smooth, projective and absolutely irreducible curve over the local field $K$.

\section{Normalized base change}\label{sec:normalized-base-change}
We want to understand the behaviour of models under base change with respect to an extension of the base field.
Usually, $L|K$ will denote a finite field extension.
Since $v_K$ is Henselian, it has a unique extension to $v_L$ to $L$, and then $(L,v_L)$ also satisfies Assumption \ref{ass:v_K}.
Given a model $\X$ of $X$, we ask ourselves whether the base change $\X\times_{\OO_K}\OO_L$ is again a model of $X_L \coloneqq X\times_K L$.
Well, almost: All properties required of a model from Definition \ref{def:model} hold except possibly normality.
We obtain a model of $X_L$, called the {\em normalized base change} of $\X$ relative to $L|K$, by normalizing $\X\times_{\OO_K}\OO_L$.\footnote{
    Here we use the fact that Assumption \ref{ass:v_K} implies that $\OO_K$ is an \emph{excellent} ring.
    It follows that $\X\times_{\OO_K}\OO_L$ is an excellent scheme, and therefore its normalization is a finite morphism.
    See \cite[Chapter 8, in particular Theorem 8.2.39 (d)]{Liu}.
}

\section{Reduced special fibres and permanence}\label{sec:reduced-special-fibres-and-permanence}
We introduce permanent models and examine their relation to reduced special fibres and weakly unramified component valuations.

We say that a model $\X$ is \textit{permanent} if $\X \times_{\OK}\OL$ is itself a model of $X$:

\begin{definition} \label{def:permanent-model}
    A model $\X$ of $X$ is called \textit{permanent}, if $\X \times_{\OK} \OL$ is normal for all finite field extensions $L|K$.
\end{definition}

We have the following theorem about permanence:

\begin{theorem} \label{thm:permanence}
    \begin{enumerate}[(i)]
        \item For a model $\X$ of $X$,  the following three statements are equivalent:
            \begin{enumerate}[(a)]
                \item $\X$ is permanent,
                \item the special fibre $\X_s$ of $\X$ is reduced,
                \item all $v\in V(\X)$ are \emph{weakly unramified} over $v_K$, i.e.\ $v$ and $v_K$ have the same value group.
            \end{enumerate}
        \item Given a model $\X$, there exists a finite extension $L|K$ such that the normalized base change of $\X$ relative to $L|K$ is permanent.
    \end{enumerate}
\end{theorem}

We will prove part (i) below; see Corollary \ref{cor:permanence-part-i}.
In Section \ref{sec:epp-and-abhyanka}, we will prove part (ii) in the special case where all multiplicities of $\X_s$
are prime to the characteristic of the residue field $k$ (Lemma \ref{lem:abhyanka}).
In later chapters, we see proofs of part (ii) in other important special cases.

\begin{remark}
    Since flatness and properness are preserved under base change, a model $\X$ of $X$ is permanent if and only if
    $\X \times_{\OK} \OL$ is a model of $X_L = X \times_K L$.
\end{remark}

\begin{proposition}[{\cite[Proposition 2.3]{ArzdorfWewers}}] \label{prop:reducedness-implies-permanence}
    If $\X_s$ is reduced, then $\X$ is permanent.
\end{proposition}

\begin{proof}
    Let $L|K$ be a finite extension and $l|k$ the extension of the corresponding residue fields.
    The special fibre of $\X' \coloneqq \X_{\OL} = \X \times_{\OK} \OL$ is $\X'_s := \X_s \times_k l$,
    which is reduced since $k$ is perfect by assumption.
    This means $\X'_s$ satisfies condition (i) from Lemma~\ref{lem:conditions-for-normality}. \\
    For condition (ii) let $\xi \in \X'_s$ be any generic point.
    The stalk $\OO_{\X', \xi}$ then is a Noetherian local domain of dimension $1$ and it holds
    \begin{align}
        \OO_{\X'_s, \xi} = \OO_{\X', \xi}/\pi_L\OO_{\X', \xi},
    \end{align}
    where $\pi_L \in \OL$ is a uniformising parameter.
    Since $\X'_s$ is reduced, $\pi_L\OO_{\X',\xi}$ is a maximal ideal, which means that
    condition (ii) from Lemma~\ref{lem:conditions-for-normality} is satisfied as well.
    Thus $\X'$ is normal and hence $\X$ is permanent.
\end{proof}

\begin{proposition}[{\cite[Corollary 2.2 (ii)]{ArzdorfWewers}}]\label{prop:criterion-reducedness}
    Let $\X$ be a model of $X$, then the following statements are equivalent.
    \begin{enumerate}[(i)]
        \item $\X_s$ is reduced.
        \item $\OO_{\X_s, \xi}$ is a field for every generic point $\xi \in \X_s$.
        \item Every $v_Z \in V(\X)$ is weakly unramified.
    \end{enumerate}
\end{proposition}

\begin{proof}
    Since $\X$ is normal, there are no embedded points in $\X_s$ by Lemma~\ref{lem:conditions-for-normality}.
    Since $\X_s$ has dimension one, it is reduced if and only if $\X_s$ is reduced in codimension zero
    which in turn is the case if $\OO_{\X_s, \xi}$ is a field for every generic point $\xi \in \X_s$. \\
    For the second equivalence recall that according to Proposition~\ref{prop:multiplicities-coincide},
    for any generic point $\xi\in \X_s$ and $Z = \overline{\lbrace\xi\rbrace}$ it holds
    \begin{align}
        [\Gamma_{v_Z}: \Gamma_K] = \length \OO_{\X_s, \xi}.
    \end{align}
    This implies that $\OO_{\X_s, \xi}$ is a field if and only if $v_Z$ is weakly unramified for all $v_Z \in V(\X)$.
\end{proof}

\begin{corollary}[Partial result] \label{cor:permanence-part-i}
    In the notation of Theorem~\ref{thm:permanence}(i), conditions (b) and (c) are equivalent, and each implies condition (a).
    The implication (a) $\Rightarrow$ (b) remains unresolved in the present text.
\end{corollary}

% AUTHOR WARNING:
\begin{authorwarning}
The implication ``permanent $\Rightarrow$ reduced special fibre'' is not proved by the preceding results and does not follow formally from the definition. A correct proof or a precise reference must be supplied.
\end{authorwarning}

\begin{proof}
    The equivalence (b) $\iff$ (c) follows from Proposition \ref{prop:criterion-reducedness}.
    The implication (b) $\Rightarrow$ (a) is Proposition \ref{prop:reducedness-implies-permanence}.
    The remaining implication (a) $\Rightarrow$ (b) is the unresolved point recorded above.
\end{proof}



\section{Extension of component valuations}\label{sec:extension-component-valuations}
To compare the special fibres before and after normalized base change, we describe how their component valuations are related by extension of valued function fields.

Let $L|K$ be a finite field extension and let $\X$ be a model of $X$.
Denote by $\X' \coloneqq (\X\times_{\OK} \OL)^\sim$ the normalized base change of $\X$ to $\OL$.
Let $\xi \in \X_s$ be a generic point and let $\xi' \in \X_s'$ be a generic point lying over $\xi$.
Since $\X$ (resp. $\X'$) is normal we have that $R \coloneqq \OO_{\X, \xi}$ (resp. $S \coloneqq \OO_{\X', \xi'}$)
is a discrete valuation ring dominating $\OK$ (resp. $\OL$) by
Lemma~\ref{lem:conditions-for-normality} (\ref{lem:conditions-for-normality-ii}), that is
\begin{itemize}
    \item $R \supset \OO_K$ (resp. $S \supset \OL$), and
    \item $\m_R \cap \OK = \pi_K\OK$ (resp. $\m_S \cap \OL = \pi_L\OL$).
\end{itemize}
We arrive at the following diagram of fields.

The function field after base change is the compositum
\begin{align}
    F_{X_L}=F_X\mathbin{\cdot}L.
\end{align}

% AUTHOR WARNING:
\begin{authorwarning}
The text assumes that a single rational subfield $K(x)\subset F_X$ can be chosen such that every component valuation $v_i$ restricts to the Gauß valuation. This simultaneous choice requires proof or a precise reference. For an individual valuation such a parameter is easy to choose, but the common choice for all components is stronger. The diagram below records the intended setup only, subject to this unresolved point.
\end{authorwarning}

\begin{figure}[H]\label{fig:valuations-and-field-extensions}
    \centering
    \resizebox{0.5\textwidth}{!}{
        \begin{tikzcd}[ampersand replacement=\&]
            {}
            \& {(F_{X_L},\{w_{i,1},\dots,w_{i,j_i}\})} \\
            {(F_X,\{v_i\})} \arrow[ru, no head]
            \& {} \\
            {}
            \& {(L(x),v_{{\scriptscriptstyle G}, L})} \arrow[uu, no head]  \\
            {(K(x),v_{{\scriptscriptstyle G}, K})} \arrow[ru, no head] \arrow[uu, no head]
            \& {} \\
            {}
            \& {(L,v_L)} \arrow[uu, no head]  \\
            {(K,v_K)} \arrow[ru, no head] \arrow[uu, no head]
            \& {} \\
        \end{tikzcd}
    }
\end{figure}

Here the $v_i$ are the valuations on $F_X$ corresponding to the generic points of the components of $\X_s$.
$w_{i, j}$ represent the valuations corresponding to the irreducible component of $\X_s'$ whose generic point lies over
the generic point of $Z_i \subset \X_s$ which corresponds to $v_i$.


\section{Eliminating ramification}\label{sec:eliminating-ramification}
We now turn to the existence of a finite base extension after which normalized base change has reduced special fibre. The required extension is obtained by eliminating ramification in the component valuations.

\begin{theorem}[{\cite[Proposition 2.3]{ArzdorfWewers}}]
    Let $\X$ be a model of $X$.
    Then there is a finite extension $L|K$ such that $\X' = (\X \times_{\OK}\OL)^\sim$ has a reduced special fibre.
\end{theorem}

% AUTHOR WARNING:
\begin{authorwarning}
A precise version of Epp's theorem appropriate to the finitely many component valuations must be stated here. In particular, the proof must specify the valued fields to which Epp is applied and the required hypotheses on their ranks and residue fields. It must also justify the passage from one component valuation to finitely many of them, explain why a single finite extension $L/K$ works simultaneously, and show that every extension $w_{i,j}$ becomes weakly unramified over $v_L$.
\end{authorwarning}

\begin{proof}
    By~\cite[Theorem 1.0]{Epp} there exists a field extension $L|K$ such that $w_{i,j}$ is weakly unramified over $v_L$.
    According to Proposition~\ref{prop:criterion-reducedness}, the special fibre of $\X'$ is reduced.
\end{proof}


\section{Epp's theorem and Abhyanka's lemma}\label{sec:epp-and-abhyanka}
The preceding existence result invokes Epp's theorem. In the tamely ramified case, the corresponding effect of a base extension on ramification indices is expressed by Abhyanka's lemma.

\begin{lemma}[Abhyanka]\label{lem:abhyanka}
    Let $L|K$ be a finite extension.
    Let $\lbrace v_1, \dots, v_n \rbrace = V(X)$ denote the extended valuation on the function field $F_X|K$.
    We have that $F_{X_L}=F_X\mathbin{\cdot}L$ and with the extended valuation denoted as $w_{1,1}, \dots, w_{n, j_n}$ respectively.
    Assume that $v_i|v_{{\scriptscriptstyle G}, K}$ are tamely ramified for all $i$.
    If $e_i \coloneqq e(v_i|v_{{\scriptscriptstyle G},K})$ divides
    $e(v_{{\scriptscriptstyle G}, L}, v_{{\scriptscriptstyle G}, K}) = e(v_L|v_K)$ for all $i$, then $w_{i,j} \in V(X_L)$
    is weakly unramified over $v_{{\scriptscriptstyle G}, L}$ for all $i,j$.
\end{lemma}

% AUTHOR WARNING:
\begin{authorwarning}
Both the statement and the proof of this lemma require revision. The notation $V(X)$ should presumably be $V(\X)$, the identity $F_{X_L}=F_X\mathbin{\cdot}L$ must be interpreted in a specified common overfield, and the extensions $w_{i,j}$ must be defined. The assertion $w_{i,j}\in V(X_L)$ uses undefined or inconsistent notation, and the ramification indices rely on the unresolved common choice of $K(x)$ recorded above. The present proof is invalid: the reduction from divisibility of ramification indices to equality is unexplained, the displayed relation between uniformisers is incorrect, and $[\Gamma_w:\Gamma_v]=1$ does not follow merely from an equality of compositum fields.
\end{authorwarning}

\begin{proof}
A correct proof must be supplied; see the warning above.
\end{proof}

\begin{comment}
\begin{proof}[Provisional argument retained for the authors]
    We will use the equivalence of Proposition~\ref{prop:criterion-reducedness} without further mention,
    that is the equivalence of the special fibre being reduced and the valuations being weakly unramified.
    Without loss of generality we may assume the following assertions:
    \begin{itemize}
        \item $n = 1$ and $j_1 = 1$ since the argument can be constructed inductively,
        \item $e \coloneqq e_1 = e(v_1|v_{{\scriptscriptstyle G},K}) = e(v_{{\scriptscriptstyle G},L}|v_{{\scriptscriptstyle G},K})$
            since we assume the residue field to be perfect, and
        \item $K$ contains the $e$-th roots of unity, i.e. $\pi_K = \pi_L^e$, since we can always replace $K$ by a suitable
            Kummer extension without violating the other assumptions and obtain a compatible diagram regarding the valuations.
    \end{itemize}
    Since $F_{X_L}=F_X\mathbin{\cdot}L$ we have that $[\Gamma_{w_1}: \Gamma_{v_1}] = 1$.
    By multiplicativity of the ramification indices we conclude $[\Gamma_{w_1}: \Gamma_{v_L}] = 1$.
\end{proof}
\end{comment}


\section{Tame and wild examples}\label{sec:tame-and-wild-examples}
The following examples illustrate how a suitable finite extension and a change of coordinates can produce a model with reduced special fibre, and why the choice of extension matters in the wildly ramified case.

\begin{example}[Obtaining a reduced special fibre after base change]
    Let $K = \QQ_3$ and consider the projective curve
    \begin{align}
        X: \, y^3 = 3x^2 + 1
    \end{align}
    over $K$.
    The equation defines a model $\X$ over $\OK=\ZZ_3$ (cf. Example \ref{ex:model-of-plane-curve}).
    The reduction $\X_s/\FF_3$ is given by $0 = y^3 - 1 = (y - 1)^3$ which means $\X_s$ is not reduced. \\
    The irreducible component is given by $Z: \, y-1$; denote the corresponding valuation by $v$.
    For $z \coloneqq y - 1$ we have
    \begin{align*}
        3x^2 &= y^3 - 1 \\
            &= (z+1)^3 - 1 \\
            &= z^3 + 3z^2 + 3z,
    \end{align*}
    where the right hand side is an Eisenstein polynomial.
    This gives
    \begin{align}
        v(z)=\frac{1}{3}v_{{\scriptscriptstyle G},K}(3x^2)=\frac{1}{3}\notin\ZZ=\Gamma_K.
    \end{align}
    Let us extend the ground field to $L = \QQ_3[\theta]$ where $\theta^3 - 3 = 0$ and apply the coordinate transformation
    $(x,y) \mapsto (x_1, 1+ \theta y_1)$.
    We then obtain a birational curve
    \begin{align}
        X_1: \, y_1^3 + \theta^2 y_1^2 + \theta y_1 = x_1^2.
    \end{align}
    Since $y_1^3 \equiv x_1^2$ (mod $\theta$), $\X_s$ is reduced.
    If $w$ denotes the extension of $v$ associated with this component after base change, then
    $w(\theta y_1)=w(y-1)=\frac{1}{3}=v_L(\theta)\in\Gamma_L$, and hence $w(y_1)=0$.
\end{example}

\begin{example}[Reduction after wildly ramified field extension]
    Consider the curve
    \begin{align}
        X: \, y^2 = 8x^3 - 2
    \end{align}
    defined over the field $K \coloneqq \QQ_2$.
    This example illustrates that the right choice of the field extension is crucial for having a reduced special fibre,
    since a wild ramified field extension is not unique.
    \begin{enumerate}[(i)]
        \item First let us consider the extension $L_1 = K(\theta_1)$ where $\theta_1^2 - 2 = 0$ and the coordinate transformation
            $(x,y) \mapsto (x_1, 2y_1 - \theta_1 - 2)$.
            We obtain a birational curve
            \begin{align}
                X_1: \, y_1^2 - (\theta_1 + 2)y_1 = 2x_1^3 - (2 + \theta_1)
            \end{align}
            The Newton polygon of the resulting quadratic in $y_1$ gives $v(y_1)=\frac14$ for an extension $v$ of
            $v_{{\scriptscriptstyle G},L_1}$ to $F_{X_1}$.
            Since $\Gamma_{L_1}=\frac12\ZZ$ and $[F_{X_1}:L_1(x_1)]=2$, it follows that $[\Gamma_{F_{X_1}}: \Gamma_{L_1}] = 2 \neq 1$.
        \item Let us consider the field extension $L_2 = K(\theta_2)$ where $\theta_2^2 + 2 = 0$
            and the coordinate transformation $(x,y) \mapsto (x_2, 2\theta_2 y_2 + \theta_2)$.
            We obtain a birational curve
            \begin{align}
                X_2: \, y_2^2 + y_2 = -x_2^3
            \end{align}
            which is given by an equation of AS-form, i.e.\ its special fibre is a reduced curve.
            Thus $[\Gamma_{F_{X_2}}: \Gamma_{L_2}] = 1$.
    \end{enumerate}
\end{example}
